\section{Beginnings}

A neural network is composed of layers which are composed of neurons. 
This structure (with a non-linear function) is very useful as it allows to approximate anything 
and shape it to any behavior we want. It's basically linear regression but on steroids. 
We will forward pass inputs and receive an output. With the output, we will calculate the loss and
run back propagation for the individual gradients and apply gradient descent on the parameters to reduce the loss value.
In turn, this will result in a model that is better suited toward the result. Repeating this process will
eventually lead the loss function reaching a minimum.

Let us define all of the internal structures needed towards a Multi Layer Perceptron (MLP).

\begin{definition}[Neuron]
    A neuron or perceptron is denoted by the following:
    \begin{equation}
        \text{N} = \sum_{i=0}^{n} w_i x_i + b , \quad n \text{ being the input size.}
    \end{equation}

    \begin{equation}
        \text{N} = 
         \begin{bmatrix}
            x_0 & x_1 & \cdots & x_i 
        \end{bmatrix}
        \begin{bmatrix}
            w_0 \\
            w_1 \\
            \vdots \\
            w_i
        \end{bmatrix}
            +
        \begin{bmatrix}
            b
        \end{bmatrix}, \quad i \text{ being the number of inputs.}
    \end{equation}
\end{definition}

\begin{definition}[Layer]
    A layer of neurons or a single layer perceptron is denoted by the following:
    \begin{equation}
        \text{L} = 
        \begin{bmatrix}
        N_0 \\
        N_1 \\
        \vdots \\
        N_i    
        \end{bmatrix}, \quad i \text{ being the amount of neurons.}
    \end{equation}
\end{definition}

\begin{definition}[MLP]
    A multi layer perceptron has multiple layers and it is denoted by the following:
    \begin{equation}
        \text{MLP} = 
        \begin{bmatrix}
            L_0 & L_1 & \cdots & L_i
        \end{bmatrix}
        =
        \begin{bmatrix}
           \begin{bmatrix}
                N_0 \\
                N_1 \\
                \vdots \\
                N_a    
          \end{bmatrix} 
          &
          \begin{bmatrix}
                N_0 \\
                N_1 \\
                \vdots \\
                N_b    
          \end{bmatrix}
          &
          \cdots
          & 
          \begin{bmatrix}
                N_0 \\
                N_1 \\
                \vdots \\
                N_c    
          \end{bmatrix}
        \end{bmatrix}, \quad a,b,c \text{ being the size of each layer.}
    \end{equation}
\end{definition}

Now that we have some infrastructure to work with, we have to show how 
forward pass would work. Imagine we have a 2 layer (1 input/middle, 1 output). 
For the input layer we shall take in two inputs and output 4. The output layer must take 
4 and we will let it output 1.
    
This is how the math would work out.

\begin{equation*}
    X_0 = 
        \begin{bmatrix}
            x_0 \\
            x_1
        \end{bmatrix},
     L_0 = 
        \begin{bmatrix}
            N_0 \\
            N_1 \\
            N_2 \\
            N_3
        \end{bmatrix}, 
    N_i = 
        \begin{bmatrix}
            x_0 & x_1 
        \end{bmatrix}
        \begin{bmatrix}
            w_0 \\
            w_1
        \end{bmatrix} + 
        \begin{bmatrix}
            b
        \end{bmatrix}
\end{equation*}

\begin{equation*}
    L_1 = 
        \begin{bmatrix}
            N_0
        \end{bmatrix},
    N_0 = 
        \begin{bmatrix}
            N_0 \\
            N_1 \\
            N_2 \\
            N_3
        \end{bmatrix}^{\top}
        \begin{bmatrix}
            w_0 \\
            w_1 \\
            w_2 \\
            w_3
        \end{bmatrix}
        =
        \begin{bmatrix}
            N_0 & N_1 & N_2 & N_3
        \end{bmatrix}
        \begin{bmatrix}
            w_0 \\
            w_1 \\
            w_2 \\
            w_3
        \end{bmatrix} + 
        \begin{bmatrix}
            b
        \end{bmatrix}
\end{equation*}